\documentclass[11pt,a4paper]{article}

\usepackage[T1]{fontenc}
\usepackage[utf8]{inputenc}
\usepackage[czech,english]{babel}
\selectlanguage{czech}
\usepackage{amsmath,amssymb,amsthm,mathtools}
\usepackage{bm}
\usepackage[a4paper,margin=2.7cm]{geometry}
\usepackage{microtype}
\usepackage{setspace}
\onehalfspacing
\usepackage{booktabs}
\usepackage{array}
\usepackage{longtable}
\usepackage{graphicx}
\usepackage{float}
\usepackage{xcolor}
\usepackage{caption}
\usepackage{listings}
\usepackage[colorlinks=true,linkcolor=blue!60!black,citecolor=blue!60!black,urlcolor=blue!60!black]{hyperref}
\newtheorem{theorem}{Věta}[section]
\newtheorem{proposition}[theorem]{Tvrzení}
\newtheorem{lemma}[theorem]{Lemma}
\newtheorem{corollary}[theorem]{Důsledek}
\theoremstyle{definition}
\newtheorem{definition}[theorem]{Definice}
\newtheorem{example}[theorem]{Příklad}
\newtheorem{remark}[theorem]{Poznámka}
\newcommand{\E}{\mathbb{E}}
\newcommand{\Var}{\mathrm{Var}}
\newcommand{\Cov}{\mathrm{Cov}}
\newcommand{\Prob}{\mathbb{P}}
\newcommand{\R}{\mathbb{R}}
\lstdefinestyle{py}{language=Python,basicstyle=\ttfamily\small,keywordstyle=\color{blue!70!black},commentstyle=\color{green!40!black},stringstyle=\color{red!50!black},showstringspaces=false,breaklines=true,frame=single,rulecolor=\color{black!20}}

\title{\vspace{-1cm}\textbf{Modelování míry storen v životním pojištění}\\\large Poissonův GLM s časovou validací a finálním holdoutem}
\author{Jaroslav Drobek}
\date{srpen 2026}

\begin{document}
\maketitle

\begin{abstract}
Tato případová studie vytváří interpretovatelný aktuárský model lapse chování v životním pojištění na datech studie Society of Actuaries 2014 Post-Level Term Lapse and Mortality Study. Počty storen jsou modelovány Poissonovým zobecněným lineárním modelem s expozicí v počtu pojistek jako offsetem. Vývoj modelu používá expanding-window časovou validaci a poslední studijní rok ponechává jako nedotčený holdout. GLM snižuje průměrnou Poissonovu devianci proti baseline tvořené portfoliovou mírou o 66,7\% v časové validaci a o 79,4\% na finálním holdoutu 2011--2012. Dominantními efekty jsou trvání smlouvy a zvýšení pojistného po skončení level-premium období. Hlavní slabinou není relativní predikční výkon, ale postupný drift absolutní kalibrace směrem k podpredikci v pozdějších letech.
\end{abstract}

\tableofcontents
\newpage

\section{Zadání úlohy}
Modelování lapse rate v životním pojištění zkoumá, jak se intenzita zániku či storna smluv mění podle produktových, klientských a smluvních charakteristik. U post-level term produktů je problém obzvlášť výrazný, protože po skončení garantovaného období konstantního pojistného může následovat prudký nárůst pojistného a s ním spojený lapse shock.

Cílem projektu je vytvořit transparentní predikční model, který:
\begin{enumerate}
\item modeluje intenzitu storna se správným zohledněním expozice,
\item zůstává interpretovatelný pomocí multiplikativních rate relativities,
\item prokazuje kvalitu na budoucích studijních letech,
\item odděluje relativní predikční výkon od absolutní kalibrace,
\item poskytuje přirozený základ pro monitoring a případnou rekalibraci.
\end{enumerate}

Veřejný repozitář projektu je dostupný na \url{https://github.com/drobekj/life-insurance-lapse-rate-modelling}.

\section{Data}
\subsection{Zdroj a struktura}
Analýza používá lapse část studie \textit{Society of Actuaries 2014 Post-Level Term Lapse and Mortality Study}. Oficiální stránka studie je \url{https://www.soa.org/resources/experience-studies/2014/research-2014-post-level-shock/}.

Lokální modelovací dataset byl rekonstruován z Excel Pivot Cache obsažené ve zdrojových materiálech studie. Každý řádek představuje agregovanou rizikovou buňku, nikoli jednotlivou pojistnou smlouvu.

\begin{center}
\begin{tabular}{lr}
\toprule
Veličina & Hodnota \\
\midrule
Zdrojové rizikové buňky & 345 627 \\
Počet proměnných & 17 \\
Modelovací buňky s kladnou expozicí & 344 725 \\
Expozice v počtu pojistek & přibližně 6,73 mil. pojistných let \\
Pozorovaná storna & přibližně 1,01 mil. \\
\bottomrule
\end{tabular}
\end{center}

Rekonstruovaný CSV soubor není ve veřejném repozitáři redistribuován. Repozitář místo toho dokumentuje oficiální zdroj a očekávaný lokální název souboru potřebný pro reprodukci notebooku.

\subsection{Vývojové období a holdout}
Časový návrh je explicitní:
\begin{itemize}
\item vývojové období: studijní roky 2000--2001 až 2010--2011,
\item finální nedotčený holdout: studijní rok 2011--2012,
\item časová validace: osm expanding-window foldů, z nichž každý predikuje jeden budoucí rok.
\end{itemize}

Toto uspořádání brání úniku informací z budoucích let do dřívějších validačních foldů a je proto pro portfolio s časově proměnlivou úrovní lapse rate informativnější než náhodný train/test split.

\section{Explorační analýza}
\subsection{Efekt trvání smlouvy}
Nejsilnějším portfoliovým vzorem je lapse shock kolem konce level-premium období. Ve vývojovém vzorku roste lapse rate z přibližně 6,7\% v trvání 6--9 na 58,3\% v trvání 10 a v následujících trváních opět klesá.

\begin{figure}[H]
\centering
\includegraphics[width=0.78\textwidth]{figures/01_lapse_rate_by_duration.pdf}
\caption{Lapse rate podle trvání smlouvy ve vývojovém vzorku}
\end{figure}

\begin{center}
\begin{tabular}{lrr}
\toprule
Trvání & Expozice & Lapse rate \\
\midrule
6--9 & 4 399 410 & 6,65\% \\
10 & 750 728 & 58,26\% \\
11 & 281 546 & 29,60\% \\
12 & 176 855 & 11,24\% \\
13+ & 442 915 & 7,00\% \\
\bottomrule
\end{tabular}
\end{center}

\subsection{Efekt zvýšení pojistného v trvání 10}
Dalším hlavním driverem je zvýšení pojistného mezi trváním 10 a 11. V trvání 10 je pozorovaná lapse rate přibližně 16,3\% pro poměr pojistného mezi 1,01 a 2,00, přesahuje 50\% při poměru nad 3 a u řady vyšších pásem překračuje 80\%.

\begin{figure}[H]
\centering
\includegraphics[width=\textwidth]{figures/02_duration10_lapse_rate_by_premium_jump.pdf}
\caption{Pozorovaná lapse rate v trvání 10 podle pásma zvýšení pojistného}
\end{figure}

V koncové části distribuce premium jump jsou expozice relativně malé, proto velmi vysoké empirické míry v nejvyšších pásmech nelze interpretovat jako dokonale monotónní strukturální vztah. Při přenosu EDA do modelové specifikace je nutné sledovat velikost expozice i stabilitu kategorií.

\section{Modelová specifikace}
\subsection{Poissonův GLM s offsetem expozice}
Pro rizikovou buňku $i$ označme $L_i$ pozorovaný počet storen a $E_i$ expozici v počtu pojistek. Model má tvar
\[
L_i \mid X_i \sim \operatorname{Poisson}(E_i\lambda_i),
\]
s logaritmickou link funkcí a offsetem expozice
\[
\log \E[L_i\mid X_i] = \log E_i + \beta_0 + X_i^\top\beta.
\]

Expozice je tedy pevně daným offsetem, nikoli volně odhadovaným koeficientem. Model tím přirozeně cílí na lapse rate a současně váží rizikové buňky podle množství expozice, které reprezentují.

Pro kategoriální koeficient $\beta_j$ lze veličinu $\exp(\beta_j)$ interpretovat jako multiplikativní rate relativity vůči referenční kategorii při ostatních faktorech konstantních.

\subsection{Finální prediktory}
Finální specifikace používá osm kategoriálních prediktorů:
\begin{itemize}
\item \texttt{DURATION},
\item \texttt{GENDER},
\item \texttt{ISSUE\_AGE\_GROUP},
\item \texttt{FACE\_AMOUNT\_BAND},
\item \texttt{POST\_LEVEL\_PREMIUM\_STRUCTURE},
\item \texttt{PREM\_JUMP\_D11\_D10},
\item \texttt{RISK\_CLASS},
\item \texttt{PREMIUM\_MODE}.
\end{itemize}

Proměnná \texttt{ISSUE\_YEAR\_GROUP} byla z finální specifikace záměrně odstraněna. Nové issue cohorts se objevují v čase a mohou se poprvé vyskytnout s významnou expozicí až v budoucích validačních letech. Zachování této proměnné by tedy činilo model zbytečně závislým na kategoriích, které ve starších trénovacích oknech strukturálně neexistují.

Kategorie jsou one-hot encoded s explicitními aktuárskými referenčními kategoriemi. Zbylé unseen levels jsou ošetřeny pomocí \texttt{handle\_unknown="ignore"}; takové případy je však vhodné monitorovat, protože neznámá úroveň je v kódování reprezentována nulami a pro daný prediktor se tak efektivně chová jako referenční vzor.

\section{Časová validace}
\subsection{Expanding-window návrh}
Validace probíhá na osmi po sobě jdoucích budoucích letech 2003--2004 až 2010--2011. V každém foldu obsahuje trénovací vzorek pouze roky předcházející validačnímu roku. Benchmarkem je jednoduchý model s konstantní portfoliovou lapse rate odhadnutou z odpovídajícího trénovacího okna.

Výkon je měřen pomocí průměrné Poissonovy deviace. Relativní zlepšení je definováno jako
\[
\text{Zlepšení} = 1 - \frac{D_{\mathrm{GLM}}}{D_{\mathrm{baseline}}}.
\]
Všech osm GLM fitů konvergovalo.

\begin{table}[H]
\centering
\small
\caption{Výsledky expanding-window validace}
\begin{tabular}{lrrr}
\toprule
Validační rok & Baseline deviance & GLM deviance & Zlepšení \\
\midrule
2003--2004 & 2,053 & 1,062 & 48,3\% \\
2004--2005 & 2,312 & 1,130 & 51,1\% \\
2005--2006 & 3,002 & 1,264 & 57,9\% \\
2006--2007 & 5,049 & 1,419 & 71,9\% \\
2007--2008 & 5,387 & 1,413 & 73,8\% \\
2008--2009 & 5,783 & 1,338 & 76,9\% \\
2009--2010 & 5,159 & 1,231 & 76,1\% \\
2010--2011 & 4,939 & 1,117 & 77,4\% \\
\midrule
Průměr & 4,210 & 1,247 & 66,7\% \\
\bottomrule
\end{tabular}
\end{table}

\begin{figure}[H]
\centering
\includegraphics[width=0.92\textwidth]{figures/03_temporal_validation_deviance.pdf}
\caption{Poissonova deviance v jednotlivých out-of-time validačních foldech}
\end{figure}

Podstatným výsledkem není pouze průměrné vítězství GLM, ale skutečnost, že model překonává baseline v každém budoucím validačním roce. Relativní predikční hodnota modelu je tedy stabilní i při výrazně odlišných portfoliových úrovních lapse rate.

\section{Kalibrace}
Predikční rozlišení a kalibrace jsou příbuzné, ale odlišné vlastnosti. Model může správně rozlišovat relativní lapse riziko mezi buňkami a současně postupně míjet agregovanou úroveň portfolia.

V souhrnu všech out-of-fold predikcí je skutečná lapse rate 14,63\%, zatímco GLM predikuje 14,11\%. Poměr Actual/Predicted je 1,037 a agregovaná kalibrační chyba činí +0,52 procentního bodu.

\begin{figure}[H]
\centering
\includegraphics[width=0.92\textwidth]{figures/04_temporal_validation_calibration.pdf}
\caption{Skutečná, GLM predikovaná a baseline lapse rate podle validačního roku}
\end{figure}

Charakter chyby se v čase mění. První validační roky jsou mírně nadhodnocené, pozdější roky se naopak posouvají k podpredikci. Například v roce 2010--2011 je skutečná portfoliová lapse rate 19,48\%, zatímco GLM predikuje 17,80\%. Baseline je výrazně níže na 13,47\%.

Toto rozlišení je prakticky důležité. Stabilní zlepšení deviace doprovázené posunem celkové úrovně naznačuje, že před úplným redevelopmentem modelu lze nejprve uvažovat jednodušší rekalibraci interceptu.

\section{Interpretace modelu}
\subsection{Relativity podle trvání}
Trvání smlouvy je dominantním fitovaným driverem. Vůči referenční skupině trvání 6--9 vycházejí následující rate relativities:

\begin{center}
\begin{tabular}{lrr}
\toprule
Efekt & Koeficient & Rate relativity \\
\midrule
Trvání 10 & 2,162 & 8,69$\times$ \\
Trvání 11 & 1,749 & 5,75$\times$ \\
Trvání 12 & 0,908 & 2,48$\times$ \\
Trvání 13+ & 0,572 & 1,77$\times$ \\
\bottomrule
\end{tabular}
\end{center}

Fitované efekty jsou menší než hrubý poměr patrný z EDA pro trvání 10, protože GLM současně kontroluje premium jump, strukturu pojistného, risk class, issue age, premium mode a další faktory.

\subsection{Další vybrané efekty}
Proměnná premium jump zůstává silným podmíněným driverem. Vůči referenčnímu pásmu 1,01--2,00 rostou fitované relativity přibližně na 1,55--1,61$\times$ v okolí pásma 5,01--8,00 a v tenkém konci distribuce už nejsou dokonale monotónní.

Významná je také post-level struktura pojistného: \textit{Premium Jump to Other} má fitted relativity přibližně 0,56$\times$ vůči \textit{Premium Jump to ART}. Premium mode má slabší, ale viditelný efekt; měsíční placení má přibližně 0,84$\times$ a biweekly mode 0,71$\times$ vůči ročnímu placení při ostatních faktorech konstantních.

Gender, issue age, face amount a risk class přidávají další segmentaci, jejich efekty jsou však slabší než struktura trvání a premium jump.

\section{Finální holdout}
Po dokončení vývoje je preprocessing encoder i Poissonův GLM znovu fitován na všech vývojových letech a jednorázově vyhodnocen na rezervovaném holdoutu 2011--2012.

\begin{itemize}
\item vývojový fit: 298 396 rizikových buněk,
\item holdout: 46 329 rizikových buněk.
\end{itemize}

\begin{table}[H]
\centering
\caption{Výsledky finálního holdoutu 2011--2012}
\begin{tabular}{lr}
\toprule
Metrika & Hodnota \\
\midrule
Baseline Poisson deviance & 5,370 \\
GLM Poisson deviance & 1,106 \\
Zlepšení proti baseline & 79,4\% \\
Skutečná lapse rate & 21,32\% \\
GLM predikovaná lapse rate & 20,06\% \\
Actual / Predicted & 1,063 \\
Kalibrační chyba & +1,26 p.b. \\
\bottomrule
\end{tabular}
\end{table}

Holdout potvrzuje centrální výsledek validace: relativní predikční výkon zůstává silný i v budoucím období. Zároveň se zvětšuje absolutní podpredikce, což odpovídá trendu kalibrace patrnému již v pozdějších validačních foldech.

\section{Praktická interpretace a monitoring}
Model je vhodný jako aktuárský portfoliový model zejména proto, že kombinuje transparentní multiplikativní strukturu s explicitním zacházením s expozicí. V produkčním prostředí by se měly pravidelně sledovat zejména:
\begin{itemize}
\item agregovaný Actual/Predicted poměr,
\item kalibrační chyba v procentních bodech,
\item Poissonova deviance proti jednoduché aktuální baseline,
\item stabilita klíčových efektů trvání a premium jump,
\item změny exposure mixu a výskyt nových kategorií.
\end{itemize}

Zhoršení deviace by znamenalo oslabování rizikové diferenciace a mohlo by vést k redevelopmentu. Naproti tomu zachování relativní deviance při téměř uniformním level shiftu přirozeně vede nejprve k rekalibraci, například úpravou interceptu.

\section{Omezení}
Studie má několik podstatných omezení:
\begin{enumerate}
\item data jsou agregované rizikové buňky, nikoli jednotlivé pojistky;
\item Poissonův model používá standardní vztah mezi střední hodnotou a variancí count procesu;
\item kategoriální efekty jsou fixní a log-aditivní, bez explicitních interakcí;
\item i při pečlivé specifikaci se mohou v budoucnu objevit nové kategorie;
\item SOA studie je historická a nemá být chápána jako náhrada současného produkčního portfolia;
\item časový drift kalibrace ukazuje, že i relativně stabilní model vyžaduje monitoring absolutní úrovně.
\end{enumerate}

Tato omezení studii neznehodnocují. Vymezují hranici mezi transparentním portfoliovým case study a produkčním modelem, který by vyžadoval aktuální firemní data, governance, monitoring a případně bohatší interakční strukturu.

\section{Závěr}
Projekt demonstruje celý aktuárský modelovací workflow:
\[
\text{rekonstrukce dat} \to \text{EDA} \to \text{Poisson GLM} \to \text{časová validace} \to \text{interpretace} \to \text{holdout}.
\]

Poissonův GLM poskytuje velmi dobrý kompromis mezi predikční kvalitou a interpretovatelností. V osmi expanding-window validačních foldech klesá průměrná deviance z 4,210 u baseline na 1,247 u GLM, tedy o 66,7\%. Na nedotčeném holdoutu 2011--2012 dosahuje zlepšení 79,4\%.

Hlavním modelovým poznatkem je dominance trvání smlouvy a premium-jump mechaniky kolem konce garantovaného level-premium období. Stejně důležitý je validační poznatek: relativní výkon zůstává silný, zatímco absolutní kalibrace postupně driftuje směrem k podpredikci. V praktickém aktuárském modelování je vhodné tyto dvě dimenze sledovat odděleně.

\newpage
\appendix
\section{Příloha: finální sada prediktorů}
\begin{lstlisting}[style=py,caption={Finální sada kategoriálních prediktorů}]
predictors = [
    "DURATION",
    "GENDER",
    "ISSUE_AGE_GROUP",
    "FACE_AMOUNT_BAND",
    "POST_LEVEL_PREMIUM_STRUCTURE",
    "PREM_JUMP_D11_D10",
    "RISK_CLASS",
    "PREMIUM_MODE",
]
\end{lstlisting}

\section{Příloha: kostra modelu}
\begin{lstlisting}[style=py,caption={Poisson GLM s offsetem expozice v počtu pojistek}]
encoder = OneHotEncoder(
    drop=[reference_categories[col] for col in predictors],
    handle_unknown="ignore",
    sparse_output=False,
)

X_train_enc = encoder.fit_transform(train[predictors])
X_valid_enc = encoder.transform(valid[predictors])

X_train_glm = sm.add_constant(X_train_enc, has_constant="add")
X_valid_glm = sm.add_constant(X_valid_enc, has_constant="add")

model = sm.GLM(
    train["LAPSE_CNT"],
    X_train_glm,
    family=sm.families.Poisson(),
    offset=np.log(train["EXPOSURE_CNT"]),
)

result = model.fit()
pred_lapses = result.predict(
    X_valid_glm,
    offset=np.log(valid["EXPOSURE_CNT"]),
)
\end{lstlisting}

\section{Příloha: reprodukovatelnost}
Veřejný repozitář obsahuje clean notebook, přesné výsledkové tabulky, grafy, requirements a dokumentaci zdrojových dat. Pro lokální reprodukci je třeba umístit \texttt{soa\_post\_level\_term\_lapse.csv} do kořene repozitáře a spustit \texttt{life\_insurance\_lapse\_rate\_modelling\_clean.ipynb} v prostředí obsahujícím pandas, NumPy, statsmodels, scikit-learn a Matplotlib.

\end{document}
